\clearpage
\addcontentsline{toc}{section}{第2セット 採点基準}

\noindent
{\Large\sffamily\bfseries 第2セット\ \ 採点基準}
\par\vspace{2mm}

\hdB{配点一覧（100点満点／200点換算併記）}
\begin{center}
\small
\renewcommand{\arraystretch}{1.3}
\begin{tabular}{@{}clcc@{}}
\toprule
\textbf{大問} & \textbf{分野・主題} & \textbf{小問の配点} & \textbf{計} \\
\midrule
1 & 数学I・A・II・III 小問集合（不定方程式・対称型三角関数・関数と数列の極限・積分方程式） & $5 + 5 + 5 + 5$ & 20 \\
2 & 数学B・III 数列の漸化式と無限級数（置き換え・階差・部分分数分解と望遠鏡和） & $4 + 5 + 5 + 6$ & 20 \\
3 & 数学C 複素数平面と一次分数変換（逆変換・円と虚軸の像・除外点・偏角の最大） & $4 + 5 + 5 + 6$ & 20 \\
4 & 数学C 空間ベクトルと球面（平面の方程式・交円上の内積・四面体の体積の最大） & $4 + 5 + 5 + 6$ & 20 \\
5 & 数学III 微分法・積分法（対数関数のグラフ・原点からの接線・面積・回転体体積） & $4 + 5 + 5 + 6$ & 20 \\
\bottomrule
\end{tabular}
\end{center}

\hdB{自己採点 成績評価基準（志望校判定・偏差値目安・到達水準）}
\par\vspace{1.0mm}\noindent
{\footnotesize 本模試（数III・複素数平面を含む理系全範囲）の自己採点得点に対応する合格目安と学習指針です。}
\par\vspace{1.5mm}
\begin{center}\footnotesize
\renewcommand{\arraystretch}{1.28}
\begin{tabular}{@{}cclp{26em}@{}}
\toprule
\textbf{得点（100点）} & \textbf{換算点（200点）} & \textbf{判定・偏差値目安} & \textbf{実力評価と今後の学習課題} \\
\midrule
$80 \sim 100$ 点 & $160 \sim 200$ 点 & \textbf{Sランク}（偏差値 68〜） & \textbf{東大理系・京大理系・国公立医学部 A判定レベル}。一次分数変換の除外点や極値・凹凸・回転体の求積まで，難度の高い数III・C記述を完璧に完遂できています。最難関実戦（第3セット）へ進む準備が整っています。 \\
$65 \sim 79$ 点 & $130 \sim 158$ 点 & \textbf{Aランク}（偏差値 60〜67） & \textbf{旧帝大理系（阪大・名大・東北・九大・北大）・難関私大理工合格水準}。大問1・2を完答し，数III微積・複素数平面の主要骨格を捉えられています。細かい減点（除外点・積分定数・増減表の漸近線）のゼロ化を目指しましょう。 \\
$45 \sim 64$ 点 & $90 \sim 128$ 点 & \textbf{Bランク}（偏差値 52〜59） & \textbf{地方国公立理系・中堅私立医学部合格水準}。数IIIの典型計算は身についていますが，空間図形や除外点の論証で減点が生じています。各大問の(1)(2)を着実に揃え，(3)の部分点を拾う実戦演習を重ねましょう。 \\
$30 \sim 44$ 点 & $60 \sim 88$ 点 & \textbf{Cランク}（偏差値 45〜51） & \textbf{理系標準記述の基礎固め段階}。無限級数の部分和極限や商の微分法・部分積分の基本公式，複素数の共役変形にまだ迷いがあります。標準問題精講や良問プラチカ等で典型解法を徹底復習しましょう。 \\
$30$ 点未満 & $60$ 点未満 & \textbf{Dランク}（偏差値 〜44） & \textbf{要・数III/Cの教科書例題徹底確認}。数III微積分の基本計算と，第1セット（数I・II・B・C基礎）の定着を最優先しましょう。 \\
\bottomrule
\end{tabular}
\end{center}

\rubricintro{無限級数における第 $N$ 部分和 $S_{N}$ の計算と $N\to\infty$ の極限操作の明記，一次分数変換における除外点 $w\neq 2$ の処理，相加・相乗平均の等号成立条件の確認など，高3秋以降の記述模試・二次試験で厳しく見られる論証要素を欠いた答案は減点の対象とする。}

% =====================================================================
% 大問1
% =====================================================================
\hdA{大問1（20点）小問集合}

\Rh{1}{1}{5}
\Seikai{個数は $4$ 個，$(x, y)=(8, 20)$ のとき最小値 $464$}
\begin{saiten}
\Sp{$5x=8(25-y)$ と互いに素な性質から一般解 $x=8k, y=25-5k$ を導いた}{+2}
\Sp{$x\geqq 1, y\geqq 1$ から $1\leqq k\leqq 4$ を導き，正整数解の個数が $4$ 個であることを求めた}{+1}
\Sp{各組（または $k$ の二次関数）の比較により，$(x, y)=(8, 20)$ で最小値 $464$ をとることを求めた}{+2}
\Ssep
\Sm{正の整数条件を見落とし，$k=0$ すなわち $(0, 25)$ を含めて $5$ 個とした}{-1}
\end{saiten}

\Rh{1}{2}{5}
\Seikai{$f(\theta)=2t^{2}+2\sqrt{2}\,t-2$（$-\sqrt{2}\leqq t\leqq\sqrt{2}$），$\theta=\dfrac{\pi}{4}$ で最大値 $6$，$\theta=\dfrac{11}{12}\pi, \dfrac{19}{12}\pi$ で最小値 $-3$}
\begin{saiten}
\Sp{$\sin 2\theta=t^{2}-1$ を用いて $f(\theta)=2t^{2}+2\sqrt{2}\,t-2$ と表した}{+1}
\Sp{合成 $t=\sqrt{2}\sin\left(\theta+\dfrac{\pi}{4}\right)$ から $-\sqrt{2}\leqq t\leqq\sqrt{2}$ を導いた}{+1}
\Sp{$t=\sqrt{2}$（$\theta=\dfrac{\pi}{4}$）のとき最大値 $6$ を求めた}{+1}
\Sp{$t=-\dfrac{\sqrt{2}}{2}$（$\theta=\dfrac{11}{12}\pi, \dfrac{19}{12}\pi$）のとき最小値 $-3$ を求めた}{+2}
\Ssep
\Sm{最小値をとる $\theta$ の値 $2$ 個のうち一方を落とした}{-1}
\end{saiten}

\Rh{1}{3}{5}
\Seikai{(i) $1$，(ii) $3$}
\begin{saiten}
\Sp{(i) 分子を有理化して $\dfrac{6n+2}{\sqrt{9n^{2}+6n+2}+3n}$ と変形した}{+1}
\Sp{(i) 分母・分子を $n$ で割って極限値 $1$ を求めた}{+1}
\Sp{(ii) $\dfrac{e^{2x}-1}{2x}\to 1$，$\dfrac{\sin 3x}{3x}\to 1$，$\dfrac{1-\cos 2x}{(2x)^{2}}\to\dfrac{1}{2}$ に分解した}{+2}
\Sp{(ii) 係数を整理して極限値 $3$ を求めた}{+1}
\Ssep
\Sm{(ii) で $\dfrac{1-\cos 2x}{(2x)^{2}}\to\dfrac{1}{2}$ の係数 $\dfrac{1}{2}$ を落として極限値を $\dfrac{3}{2}$ とした}{-1}
\end{saiten}

\Rh{1}{4}{5}
\Seikai{$f(x)=\sin x-\cos x$}
\begin{saiten}
\Sp{$\sin(x-t)=\sin x\cos t-\cos x\sin t$ で展開し，定積分を定数 $A, B$ とおいて $f(x)=(2+A)\sin x-B\cos x$ と表した}{+2}
\Sp{$A, B$ の定義式に代入して定積分を計算し，連立方程式 $A=-B, B=2+A$ を導いた}{+2}
\Sp{$A=-1, B=1$ を求めて $f(x)=\sin x-\cos x$ を決定した}{+1}
\Ssep
\Sm{積分変数 $t$ と変数 $x$ を混同し，$\sin(x-t)$ のまま積分の外に出した}{-3}
\end{saiten}

% =====================================================================
% 大問2
% =====================================================================
\hdA{大問2（20点）数列の漸化式と無限級数}

\Rh{2}{1}{4}
\Seikai{$a_{2}=\dfrac{7}{6}$，$a_{3}=\dfrac{5}{4}$，$b_{n+1}-b_{n}=3n+2$}
\begin{saiten}
\Sp{$a_{2}=\dfrac{7}{6}$，$a_{3}=\dfrac{5}{4}$ をともに正しく求めた}{+2}
\Sp{漸化式の両辺に $(n+1)(n+2)$ を掛けて $b_{n+1}-b_{n}=3n+2$ を導いた}{+2}
\Ssep
\Sm{$a_{3}$ の約分または通分を誤った}{-1}
\end{saiten}

\Rh{2}{2}{5}
\Seikai{$b_{n}=\dfrac{n(3n+1)}{2}$，$a_{n}=\dfrac{3n+1}{2(n+1)}$}
\begin{saiten}
\Sp{階差数列を用いて $b_{n}=b_{1}+\sum_{k=1}^{n-1}(3k+2)$（$n\geqq 2$）を立式した}{+2}
\Sp{和を計算し，$n=1$ の確認を含めて $b_{n}=\dfrac{n(3n+1)}{2}$ を求めた}{+2}
\Sp{$a_{n}=\dfrac{b_{n}}{n(n+1)}=\dfrac{3n+1}{2(n+1)}$ を導いた}{+1}
\Ssep
\Sm{階差の和の上端を $n$ と誤った（$\sum_{k=1}^{n}$ とした）}{-2}
\end{saiten}

\Rh{2}{3}{5}
\Seikai{$L=\dfrac{3}{2}$，収束しその和は $1$}
\begin{saiten}
\Sp{$L=\lim\limits_{n\to\infty}a_{n}=\dfrac{3}{2}$ を求めた}{+1}
\Sp{$\dfrac{L-a_{n}}{n}=\dfrac{1}{n(n+1)}=\dfrac{1}{n}-\dfrac{1}{n+1}$ と部分分数分解した}{+2}
\Sp{第 $N$ 部分和 $S_{N}=1-\dfrac{1}{N+1}$ を求め，$N\to\infty$ として収束と和 $1$ を導いた}{+2}
\Ssep
\Sm{第 $N$ 部分和 $S_{N}$ を経ずに形式的に無限項を打ち消して答えのみを出した}{-1}
\end{saiten}

\Rh{2}{4}{6}
\Seikai{$\dfrac{7}{8}$}
\begin{saiten}
\Sp{$\dfrac{a_{n}}{n(n+2)}=\dfrac{3n+1}{2n(n+1)(n+2)}$ を導いた}{+1}
\Sp{$\dfrac{1}{4}\left(\dfrac{1}{n}-\dfrac{1}{n+1}\right)+\dfrac{5}{4}\left(\dfrac{1}{n+1}-\dfrac{1}{n+2}\right)$（または同等な部分分数分解）を正しく行った}{+2}
\Sp{第 $N$ 部分和 $T_{N}=\dfrac{1}{4}\left(1-\dfrac{1}{N+1}\right)+\dfrac{5}{4}\left(\dfrac{1}{2}-\dfrac{1}{N+2}\right)$ を求めた}{+2}
\Sp{$N\to\infty$ の極限をとって和 $\dfrac{7}{8}$ を求めた}{+1}
\Ssep
\Sm{部分分数分解の係数を誤った}{-2}
\end{saiten}

% =====================================================================
% 大問3
% =====================================================================
\hdA{大問3（20点）複素数平面と一次分数変換}

\Rh{3}{1}{4}
\Seikai{$z=\dfrac{2w-1}{w-2}$，単位円 $|w|=1$}
\begin{saiten}
\Sp{$(w-2)z=2w-1$ から $w\neq 2$ を示し，$z=\dfrac{2w-1}{w-2}$ を導いた}{+2}
\Sp{$|2w-1|=|w-2|$ を展開・整理して $|w|=1$（単位円）を導いた}{+2}
\Ssep
\Sm{$w\neq 2$ の確認がない}{-1}
\end{saiten}

\Rh{3}{2}{5}
\Seikai{点 $\dfrac{5}{4}$ を中心とする半径 $\dfrac{3}{4}$ の円（ただし点 $2$ を除く）と図示}
\begin{saiten}
\Sp{虚軸上の条件 $z+\bar{z}=0$（または $|z-1|=|z+1|$ 等）に $z=\dfrac{2w-1}{w-2}$ を代入した}{+1}
\Sp{式を整理して $\left|w-\dfrac{5}{4}\right|=\dfrac{3}{4}$ を導いた}{+2}
\Sp{除外点 $w\neq 2$ を明記し，中心 $\dfrac{5}{4}$・半径 $\dfrac{3}{4}$・実軸との交点 $\dfrac{1}{2}, 2$（白丸）を正しく図示した}{+2}
\Ssep
\Sm{除外点 $w\neq 2$（点 $2$ を除くこと）を記述・図示ともに落とした}{-2}
\end{saiten}

\Rh{3}{3}{5}
\Seikai{$z=0$ のとき最小値 $\dfrac{1}{2}$，とり得る値の範囲は $\dfrac{1}{2}\leqq |w|<2$}
\begin{saiten}
\Sp{円の中心 $\dfrac{5}{4}$ と半径 $\dfrac{3}{4}$ から $|w|$ の最小値が $\dfrac{5}{4}-\dfrac{3}{4}=\dfrac{1}{2}$ であることを求めた}{+2}
\Sp{$w=\dfrac{1}{2}$ に対応する $z=0$ を求めた}{+1}
\Sp{点 $w=2$ が除外点であることから，上端に等号を含まない範囲 $\dfrac{1}{2}\leqq |w|<2$ を求めた}{+2}
\Ssep
\Sm{除外点 $w=2$ を考慮せず，$\dfrac{1}{2}\leqq |w|\leqq 2$ と上端に等号を入れた}{-1}
\end{saiten}

\Rh{3}{4}{6}
\Seikai{$\tan\theta$ の最大値 $\dfrac{3}{4}$（$w=\dfrac{4+3i}{5}$，$z=-i$）}
\begin{saiten}
\Sp{$z=ti\ (t<0)$ を代入して $w=\dfrac{2(t^{2}+1)-3ti}{t^{2}+4}$ と表し，$\tan\theta=\dfrac{-3t}{2(t^{2}+1)}$ を導いた（または円の上半円への接線条件を立てた）}{+2}
\Sp{相加・相乗平均の関係（または接線の直角三角形）から $\tan\theta$ の最大値 $\dfrac{3}{4}$ を求めた}{+2}
\Sp{最大となるときの $z=-i$（$t=-1$）および $w=\dfrac{4+3i}{5}$ をともに求めた}{+2}
\Ssep
\Sm{$t<0$ の符号処理を誤り，$\tan\theta=-\dfrac{3}{4}$ とした}{-2}
\end{saiten}

% =====================================================================
% 大問4
% =====================================================================
\hdA{大問4（20点）空間ベクトルと球面の切り口}

\Rh{4}{1}{4}
\Seikai{$\alpha: x+y+z=3$，$\Pt{H}(1, 1, 1)$，$\Pt{OH}=\sqrt{3}$}
\begin{saiten}
\Sp{$3$ 点 $\Pt{A}, \Pt{B}, \Pt{C}$ を通ることから平面 $\alpha: x+y+z=3$ を求めた}{+1}
\Sp{法線ベクトル $(1,1,1)$（または対称性・直交条件）から $\Pt{H}(1, 1, 1)$ を求めた}{+2}
\Sp{$\Pt{OH}=\sqrt{3}$ を求めた}{+1}
\Ssep
\Sm{$\Pt{H}(1,1,1)$ の根拠（$\Vec{OH}\perp\alpha$）の記述がない}{-1}
\end{saiten}

\Rh{4}{2}{5}
\Seikai{$r=\sqrt{6}$，最大値 $9$，最小値 $-3$}
\begin{saiten}
\Sp{三平方の定理 $r=\sqrt{3^{2}-(\sqrt{3})^{2}}$ から $r=\sqrt{6}$ を求めた}{+1}
\Sp{$\Vec{OP}=\Vec{OH}+\Vec{HP}$ と分解し，$\Vec{OA}\cdot\Vec{OP}=3+\Vec{HA}\cdot\Vec{HP}=3+6\cos\varphi$ を導いた}{+2}
\Sp{最大値 $9$ と最小値 $-3$ をともに正しく求めた}{+2}
\Ssep
\Sm{点 $\Pt{P}$ が球面 $S$ 全体を動くものとして最大値 $9$・最小値 $-9$ とした}{-2}
\end{saiten}

\Rh{4}{3}{5}
\Seikai{$\triangle\Pt{BCD}=\dfrac{3\sqrt{3}}{2}$}
\begin{saiten}
\Sp{点 $\Pt{D}(-1, 2, 2)$ が平面 $\alpha$ の式と球面 $S$ の式をともに満たすことを確認した}{+1}
\Sp{$\Vec{BC}=(0, -3, 3)$，$\Vec{BD}=(-1, -1, 2)$ の大きさと内積を計算した}{+2}
\Sp{面積公式から $\triangle\Pt{BCD}=\dfrac{3\sqrt{3}}{2}$ を求めた}{+2}
\Ssep
\Sm{点 $\Pt{D}\in K$ の確認を省いた}{-1}
\end{saiten}

\Rh{4}{4}{6}
\Seikai{体積の最大値 $\dfrac{3\sqrt{3}+3}{2}$，$\Pt{Q}(-\sqrt{3}, -\sqrt{3}, -\sqrt{3})$}
\begin{saiten}
\Sp{点 $\Pt{Q}$ から平面 $\alpha$ への距離（高さ $h$）が最大となるときに体積が最大となることを指摘した}{+1}
\Sp{高さの最大値が $h_{\max}=\Pt{OH}+R=\sqrt{3}+3$ であることを導いた}{+2}
\Sp{そのときの点 $\Pt{Q}(-\sqrt{3}, -\sqrt{3}, -\sqrt{3})$ を求めた}{+1}
\Sp{体積の最大値 $V_{\max}=\dfrac{1}{3}\cdot\dfrac{3\sqrt{3}}{2}\cdot(\sqrt{3}+3)=\dfrac{3\sqrt{3}+3}{2}$ を求めた}{+2}
\Ssep
\Sm{$\Vec{OQ}$ の向きを逆にとり，$h=3-\sqrt{3}$ または $\Pt{Q}(\sqrt{3}, \sqrt{3}, \sqrt{3})$ とした}{-2}
\end{saiten}

% =====================================================================
% 大問5
% =====================================================================
\hdA{大問5（20点）対数関数の微積分・面積・回転体の体積}

\Rh{5}{1}{4}
\Seikai{極大値 $\dfrac{1}{e}$（$x=e$），変曲点 $\left(e^{3/2}, \dfrac{3}{2e^{3/2}}\right)$，$\lim\limits_{x\to +0}f(x)=-\infty$，$\lim\limits_{x\to\infty}f(x)=0$，概形}
\begin{saiten}
\Sp{$f'(x)=\dfrac{1-\log x}{x^{2}}$，$f''(x)=\dfrac{2\log x-3}{x^{3}}$ を計算し，極大値 $\dfrac{1}{e}$（$x=e$）と変曲点 $\left(e^{3/2}, \dfrac{3}{2e^{3/2}}\right)$ を求めた}{+2}
\Sp{両端の極限 $\lim\limits_{x\to +0}f(x)=-\infty$，$\lim\limits_{x\to\infty}f(x)=0$ を求め，増減・凹凸と漸近線を反映した概形を描いた}{+2}
\Ssep
\Sm{第2次導関数 $f''(x)$ と変曲点の調査を省いた}{-2}
\end{saiten}

\Rh{5}{2}{5}
\Seikai{$\ell: y=\dfrac{1}{2e}x$，$\Pt{T}\left(\sqrt{e}, \dfrac{1}{2\sqrt{e}}\right)$，不等式の証明}
\begin{saiten}
\Sp{接点を $\Pt{T}(t, f(t))$ とおいて接線の方程式を立て，原点を通る条件から $t=\sqrt{e}$ を導いた}{+2}
\Sp{接線 $\ell: y=\dfrac{1}{2e}x$ と接点 $\Pt{T}\left(\sqrt{e}, \dfrac{1}{2\sqrt{e}}\right)$ を求めた}{+1}
\Sp{$g(x)=\dfrac{x^{2}}{2e}-\log x$（または $\dfrac{\log x}{x^{2}}$）の増減を調べて $x=\sqrt{e}$ で最小値 $0$（または最大値 $\dfrac{1}{2e}$）をとることを示し，不等式を証明した}{+2}
\Ssep
\Sm{「グラフより明らか」として微分による増減の根拠なしに不等式を結論づけた}{-2}
\end{saiten}

\Rh{5}{3}{5}
\Seikai{$S=\dfrac{1}{8}$}
\begin{saiten}
\Sp{求める面積を $S=\dfrac{1}{2}\cdot\sqrt{e}\cdot\dfrac{1}{2\sqrt{e}}-\int_{1}^{\sqrt{e}}\dfrac{\log x}{x}\,dx$（または区間を分けた定積分）として正しく立式した}{+2}
\Sp{原始関数 $\dfrac{1}{2}(\log x)^{2}$ を用いて $\int_{1}^{\sqrt{e}}\dfrac{\log x}{x}\,dx=\dfrac{1}{8}$ を計算した}{+2}
\Sp{$S=\dfrac{1}{4}-\dfrac{1}{8}=\dfrac{1}{8}$ を求めた}{+1}
\Ssep
\Sm{$0<x\leqq 1$ の部分で曲線 $C$ を積分に含めてしまった（$x$ 軸との境界 $(1,0)$ を見落とした）}{-2}
\end{saiten}

\Rh{5}{4}{6}
\Seikai{$V=\pi\left(2-\dfrac{5}{e}\right)$}
\begin{saiten}
\Sp{回転体の体積を $V=\pi\int_{1}^{e}\dfrac{(\log x)^{2}}{x^{2}}\,dx$ と立式した}{+1}
\Sp{第1回の部分積分により $\left[-\dfrac{(\log x)^{2}}{x}\right]_{1}^{e}+2\int_{1}^{e}\dfrac{\log x}{x^{2}}\,dx$ と変形した}{+2}
\Sp{第2回の部分積分により $\int_{1}^{e}\dfrac{\log x}{x^{2}}\,dx=1-\dfrac{2}{e}$ を求めた}{+2}
\Sp{全体を整理して $V=\pi\left(2-\dfrac{5}{e}\right)$ を求めた}{+1}
\Ssep
\Sm{回転体の公式の $\pi$ を付け忘れた}{-1}
\Sm{部分積分の符号を誤った}{-2}
\end{saiten}

